\documentclass[12pt]{report}

\usepackage[margin=1in]{geometry}
\usepackage{amsmath}
\usepackage{amssymb}
\usepackage{booktabs}
\usepackage{enumitem}
\usepackage{xcolor}
\usepackage{hyperref}
\usepackage{fancyhdr}

\hypersetup{
  colorlinks=true,
  linkcolor=blue,
  urlcolor=blue
}

\setlength{\parindent}{0pt}
\setlength{\parskip}{0.55em}
\pagestyle{fancy}
\fancyhf{}
\fancyhead[L]{\footnotesize ENGS 101: Introduction to Reinforcement Learning}
\fancyhead[R]{\footnotesize Final Project}
\fancyfoot[C]{\thepage}

\newcommand{\duedate}{TBD}
\newcommand{\repositoryurl}{https://github.com/cccat6/engs101-cyber-rl}

\begin{document}

\noindent
\fbox{
  \parbox{\dimexpr\textwidth-2\fboxsep-2\fboxrule}{
    \textbf{ENGS 101: Introduction to Reinforcement Learning}
    \hfill \textbf{Final Project}\\[0.6em]
    \begin{center}
      {\Large \textbf{Red--Blue Cybersecurity RL Challenge}}
    \end{center}
    Instructor: Prof.\ Peter Chin
    \hfill Due: \duedate
  }
}

\vspace{0.5em}

\begin{enumerate}[leftmargin=0.25in,label=$\diamond$,itemsep=0.15em]
  \item Submit executable code through the course submission website.
  \item Submit one agent implementation. The same code must support both the
        Red and Blue roles.
  \item You may use the Python standard library, PyTorch, and NumPy. You may
        not use an existing reinforcement-learning implementation or framework.
  \item The course server performs the official training and evaluation.
        Local results are for development only.
\end{enumerate}

\section*{1. Objective}

Design and implement a reinforcement-learning agent for a simulated
cybersecurity environment. Red attempts to discover, compromise, and impact
network assets. Blue attempts to detect activity, deploy decoys, remove access,
and restore hosts.

Your code will be trained separately as Red and Blue. The resulting frozen
policies will play against the policies submitted by the rest of the class and
the course baselines.

\section*{2. Starter repository}

Clone the public repository:

\begin{verbatim}
git clone https://github.com/cccat6/engs101-cyber-rl
cd engs101-cyber-rl
python -m venv .venv
python -m pip install -r requirements.txt
\end{verbatim}

The repository contains a small public environment and one required student
file:

\begin{center}
\texttt{student\_agent.py}
\end{center}

The public environment is smaller than the hidden server environment. Do not
hard-code observation dimensions, action counts, host names, or a particular
topology.

\section*{3. Required interface}

Keep the following function:

\begin{verbatim}
def create_agent(config):
    return YourAgent(config)
\end{verbatim}

The returned object must expose:

\begin{verbatim}
agent.model
agent.reset(seed)
agent.act(observation, training)
agent.learn(transition)
\end{verbatim}

\texttt{agent.model} must be a \texttt{torch.nn.Module}. The other three
members have the following meanings:

\begin{center}
\begin{tabular}{p{0.30\textwidth}p{0.62\textwidth}}
\toprule
\textbf{Member} & \textbf{Purpose} \\
\midrule
\texttt{reset(seed)} &
Reset episode-local state. \\
\texttt{act(observation, training)} &
Return one valid integer action. \\
\texttt{learn(transition)} &
Store the transition and perform any training update. \\
\bottomrule
\end{tabular}
\end{center}

The runner supplies this configuration:

\begin{verbatim}
config.role
config.observation_dim
config.num_actions
config.total_steps
config.seed
config.device
\end{verbatim}

Always construct the model from the supplied dimensions and move it to
\texttt{config.device}. The server normally supplies a CUDA device; local
development may use CPU, CUDA, Intel XPU, or Apple MPS.

\section*{4. Observations and transitions}

The policy observation has two fields:

\begin{verbatim}
observation.features
observation.action_mask
\end{verbatim}

\texttt{features} is a one-dimensional \texttt{float32} NumPy array.
\texttt{action\_mask} is a Boolean array with one entry per action.
\texttt{act} must return an action whose mask entry is true.

During training, \texttt{learn} receives:

\begin{verbatim}
transition.observation
transition.action
transition.reward
transition.next_observation
transition.done
\end{verbatim}

The runner calls \texttt{learn} exactly once per environment step. Evaluation
does not call \texttt{learn}.

\section*{5. Implementation rule}

You must implement the RL algorithm yourself using PyTorch tensor operations,
neural-network modules, and optimizers. Examples of allowed components include
\texttt{torch.nn.Linear}, activation functions, automatic differentiation,
and \texttt{torch.optim}.

Do not import an implementation of DQN, PPO, actor--critic, replay buffers,
training loops, or policies from an RL package. In particular, do not use
Stable-Baselines3, RLlib, CleanRL, Tianshou, TorchRL algorithms, or equivalent
ready-made implementations.

\section*{6. Local run}

From the repository root, run:

\begin{verbatim}
python -m cyberrl_eval.submission_cli --steps 1000 --device auto
\end{verbatim}

This single command:

\begin{enumerate}[leftmargin=0.3in]
  \item trains the submitted code as Red;
  \item saves and reloads the Red model;
  \item tests the frozen Red model;
  \item repeats the same process as Blue.
\end{enumerate}

The provided starter agent is only an interface example. It selects random
valid actions during training and does not learn.

\section*{7. Server training and evaluation}

The course server supplies the official environment-step budget, random seeds,
training scenarios, device, and evaluation episodes. These values are the same
for all submissions.

After training, the server saves the weights in \texttt{agent.model}. It then
creates a fresh agent, reloads those weights, and performs frozen inference.

The public and hidden environments use the same interface but may have
different dimensions and network topologies.

\section*{8. Reward and score}

The environment supplies a per-step reward for learning. The leaderboard score
is computed separately from the mission loss over the whole episode.

For an episode of length $T$, let

\[
D = \frac{1}{T}\sum_{t=1}^{T}\text{mission\_loss}_t.
\]

The corresponding role scores are

\[
\text{RedScore}=100D,
\qquad
\text{BlueScore}=100(1-D).
\]

The leaderboard aggregates matches across the hidden evaluation scenarios.
Training rewards are not leaderboard points.

\section*{9. Leaderboard}

Each accepted submission produces one frozen Red policy and one frozen Blue
policy. The server evaluates the Red policies against the Blue policies from
all accepted submissions and the course baselines.

When a new submission is accepted, the server rebuilds the comparison matrix
using the saved frozen policies. Old policies are re-evaluated, not retrained.
The website reports leaderboard results; hidden scenarios and trajectories are
not released.

\section*{10. Submission}

Submit \texttt{student\_agent.py} through the course website. Do not submit
pretrained weights. The server trains a fresh Red model and a fresh Blue model
from your code.

Submissions must:

\begin{enumerate}[leftmargin=0.3in]
  \item import successfully in the fixed course environment;
  \item follow the required interface;
  \item return only valid actions;
  \item use only the permitted libraries;
  \item complete within the server resource limits.
\end{enumerate}

The due date, official training budget, resource limits, and grading rubric
will be provided by the instructor.

\end{document}
